% Distribución causal exacta del delta Ley 9; contenido conservado sin síntesis.
\Needspace{28\baselineskip}
\section{Bifurcation causale postérieure à \texorpdfstring{\(\alpha_{\mathrm{HMT}}\)}{alpha HMT}}
\label{sec:l9s102:bifurcacion-accion-angular}

La généalogie de l'action et du couple angulaire forme un diagramme bifurqué.
La branche régionale produit la mantisse réduite \(\eta_{\mathrm{ret}}\) ; la branche
déterminantale produit l'ordre décimal \(\varepsilon_H=-34\). Les deux
convergent dans la droite d'action. Le couple angulaire est construit dans une branche
sœur à partir de \(\alpha_{\mathrm{HMT}}\) et de
\(\eta_{\mathrm{ret}}\) :
\begin{equation}
\begin{tikzcd}[column sep=small,row sep=normal]
\alpha_{\mathrm{HMT}},\pi,\varphi
  \arrow[r] \arrow[d]
& H_5,\mathscr C_\pi
  \arrow[d]
&& \alpha_{\mathrm{HMT}},\varphi,H_4
  \arrow[d] & {}\\
x_A,D_A,\rho_\pi
  \arrow[r]
& \eta_{\mathrm{ret}}
  \arrow[r] \arrow[d]
& \hbar_{\mathrm{ret}}\in\mathcal L_S
& \Sigma_H,\varepsilon_H=-34
  \arrow[l] & {}\\
& (A_{\deg},C^*_{\deg})
  \arrow[r]
& (A_{\mathrm{rad}},C^*_{\mathrm{rad}})
  \arrow[r]
& (q_+,q_-)
  \arrow[r]
& V_{90,120}.
\end{tikzcd}
\label{diag:l9s102:dag-accion-angular}
\end{equation}
L'ordre d'exposition situe la dérivation de Planck avant la publication
angulaire complète ; le diagramme conserve simultanément la dépendance
mathématique exacte de chaque sortie.

\section{Conoyau local et dérivation de l'ordre absolu \texorpdfstring{\(10^{-34}\)}{10^-34}}
\label{sec:l9s102:orden-accion}

La forme normale de Smith du bloc dodécaphasique local est
\begin{equation}
 \operatorname{SNF}(H_4)=\operatorname{diag}(1,2,2,4),
 \qquad
 \operatorname{coker}(H_4)
 \cong\mathbb Z/2\mathbb Z\oplus\mathbb Z/2\mathbb Z
 \oplus\mathbb Z/4\mathbb Z.
 \label{eq:l9s102:snf-h4}
\end{equation}
Le conoyau local est d'ordre \(16\). Pour la section constante
\(s_\alpha(g)=\alpha_{\mathrm{HMT}}\), la norme locale et une action de
\(\varphi\) produisent
\begin{equation}
 \boxed{
 \Sigma_H=\varphi\mathcal N_{G_H}(s_\alpha)
 =\varphi\alpha_{\mathrm{HMT}}^{16}.}
 \label{eq:l9s102:sigma-h}
\end{equation}
Les inégalités exactes
\begin{equation}
 \alpha_{\mathrm{HMT}}^{16}<10^{-34}
 <\varphi\alpha_{\mathrm{HMT}}^{16}<10^{-33}
 \label{eq:l9s102:orden-34}
\end{equation}
déterminent
\begin{equation}
 \boxed{
 \varepsilon_H=\left\lfloor\log_{10}\Sigma_H\right\rfloor=-34.}
 \label{eq:l9s102:epsilon-h}
\end{equation}
